\HandoutSetup
  {NE 630}
  {01}
  {Mass, Energy, and Nuclear Reactions}
  {Read FNRP Sections 1.1--1.3. Define the key terms in your notes and use the learning objectives to guide the reading.}

\begin{document}
\MakeHandoutTitle

\begin{handoutcolumns}

\begin{fixedbox}{Macro objective}
Determine and quantify the reactants, products,
and energy inputs/outputs of nuclear reactions.
\end{fixedbox}

\TermStrip{nuclide; atomic number; atomic mass; mass number; nucleon; rest mass; total energy; kinetic energy; binary reaction; radioactive decay; $Q$ value; binding energy; energy density}

\HandoutSection{1}{Nuclides and reaction language}

\begin{boardbox}{Read the symbol}
\[
  {}^{A}_{Z}X,
  \qquad A=Z+N
\]
\begin{tabularx}{\linewidth}{@{}>{\bfseries}l X@{}}
$X$ & \Blank[1.85in]\\[1pt]
$Z$ & \Blank[1.85in]\\[1pt]
$A$ & \Blank[1.85in]\\[1pt]
$N$ & \Blank[1.85in]
\end{tabularx}
\end{boardbox}

\begin{fixedbox}{Reaction forms used in this lesson}
\[
  A+B\longrightarrow C+D,
  \qquad A(B,C)D
\]
\[
  A\longrightarrow B+C
  \qquad\text{(radioactive decay)}
\]
Organize the initial and final values before using mass data:
\[
\begin{array}{c|cc}
 & \text{initial} & \text{final}\\ \hline
\text{mass number} & \Blank[0.46in] & \Blank[0.46in]\\
\text{atomic number} & \Blank[0.46in] & \Blank[0.46in]
\end{array}
\]
\end{fixedbox}

% Instructor cue: ask for syllabus questions, introduce the text and homework
% process, and use the notation table to reactivate NE 495 vocabulary.

\HandoutSection{2}{Mass change and $Q$ values}

\begin{fixedbox}{Relations and conversion}
\[
  \Delta E=\Delta m c^2,
  \qquad
  Q=\left(\sum m_{\mathrm{initial}}-\sum m_{\mathrm{final}}\right)c^2
\]
\[
  1\unit{u}=931.5\unit{MeV}/c^2
\]
The reaction rearranges bound nucleons (or electrons in a chemical
reaction). Nuclear changes are typically on the $\unit{MeV}$ scale;
chemical changes are typically on the $\unit{eV}$ scale.
\textbf{Board check:} If $Q>0$, which mass sum is larger, and what does
that imply for a decay?
\end{fixedbox}

\switchcolumn

\HandoutSection{3}{Relating neutron energies}

\begin{fixedbox}{Exact and classical expressions}
\[
  E=mc^2,
  \qquad
  m=\frac{m_0}{\sqrt{1-(v/c)^2}},
  \qquad \beta=\frac{v}{c}
\]
\[
  \Ktrue=m_0c^2\left(\frac{1}{\sqrt{1-\beta^2}}-1\right),
  \qquad
  \Kclassical=\frac12m_0v^2
\]
\[
  \%\,\mathrm{err}=100\frac{\Kclassical-\Ktrue}{\Ktrue}
\]
\textbf{Neutron data:}
$m_0=1.008664\unit{u}$,
$m_0c^2\approx 939.57\unit{MeV}$,
$c=2.9979\times10^8\unit{m/s}$.
\end{fixedbox}

\begin{boardbox}{Speed associated with a given $\Ktrue$}
Starting from the exact kinetic energy, isolate $\beta^2$:
\[
  \beta^2
  =1-\left(\frac{m_0c^2}{\Ktrue+m_0c^2}\right)^2
  =\frac{v^2}{c^2}.
\]
Use this result to express $\Kclassical$ entirely in terms of $\Ktrue$
and $m_0c^2$.
\RuledLines{2}
\end{boardbox}

\begin{checkpoint}[Threshold question]
Find the maximum true neutron kinetic energy for which the magnitude of
the classical error is no more than $1\%$.
\RuledLines{2}
\textbf{Board result:}
$\Ktrue\approx\RevealBlank[0.78in]{6.3 MeV}$.
\end{checkpoint}

\begin{boardbox}{Homework bridge: a $2.0\unit{MeV}$ fission neutron}
\begin{tabularx}{\linewidth}{@{}l X l X@{}}
$\beta^2$ & \Blank[0.72in] & $v$ & \Blank[0.72in]\\[2pt]
$\Kclassical$ & \Blank[0.72in] & error & \Blank[0.72in]
\end{tabularx}
\end{boardbox}

\begin{takeawaybox}[Interpret the approximation]
The classical expression \RevealBlank[1.05in]{underestimates} the exact
kinetic energy. Connect the $1\%$ threshold to the lesson objective that
neutron kinematics below about $10\unit{MeV}$ can often be treated classically.
\end{takeawaybox}

% Instructor cue: ask what each term in E=mc^2 means before the algebra;
% after obtaining the error, emphasize that its sign is negative.

\end{handoutcolumns}

\newpage
\begin{handoutcolumns}

\HandoutSection{4}{Reaction completion and $Q$-value practice}

\begin{boardbox}{Complete each reaction; then compute $Q$}
Use one reaction as the board example; retain the others for practice.
\begin{enumerate}[label=(\alph*)]
\item ${}^{9}_{\Blank[0.24in]}\mathrm{Be}+{}^{4}_{2}\mathrm{He}
      \longrightarrow \Blank[0.88in]+{}^{1}_{1}\mathrm{H}$
\item ${}^{60}_{\Blank[0.24in]}\mathrm{Co}
      \longrightarrow \Blank[0.92in]+{}^{0}_{-1}\mathrm{e}$
\item ${}^{7}_{3}\mathrm{Li}+{}^{1}_{1}\mathrm{H}
      \longrightarrow \Blank[0.88in]+{}^{4}_{2}\mathrm{He}$
\item ${}^{10}_{5}\mathrm{B}+{}^{4}_{2}\mathrm{He}
      \longrightarrow \Blank[0.88in]+{}^{1}_{1}\mathrm{H}$
\end{enumerate}
\begin{tabularx}{\linewidth}{@{}>{\bfseries}l X >{\bfseries}l X@{}}
$Q_a$ & \Blank[0.74in] & $Q_b$ & \Blank[0.74in]\\[1pt]
$Q_c$ & \Blank[0.74in] & $Q_d$ & \Blank[0.74in]
\end{tabularx}
\end{boardbox}

\begin{fixedbox}{Atomic-mass data used in the assigned problems ($\unit{u}$)}
\centering\scriptsize
\begin{tabular}{@{}lr@{\quad}lr@{}}
\toprule
nuclide & mass & nuclide & mass\\
\midrule
${}^{1}\mathrm{H}$   & 1.007825031898 & ${}^{4}\mathrm{He}$ & 4.00260325413\\
${}^{7}\mathrm{Li}$  & 7.01600343426  & ${}^{9}\mathrm{Be}$ & 9.012183062\\
${}^{10}\mathrm{B}$  & 10.012936862   & ${}^{12}\mathrm{B}$ & 12.014352638\\
${}^{13}\mathrm{C}$  & 13.00335483534 & ${}^{60}\mathrm{Co}$& 59.933815536\\
${}^{60}\mathrm{Ni}$ & 59.930785129   & $n$                   & 1.008664\\
\bottomrule
\end{tabular}
\end{fixedbox}

\begin{checkpoint}[Atomic masses and decay]
Before evaluating the $\beta^-$-decay $Q$ value with the supplied
\emph{atomic} masses, show which electron masses are already contained in
the parent and daughter entries.
\RuledLines{1}
\end{checkpoint}

\begin{boardbox}{Mass defect and binding energy per nucleon}
Reconstruct the definitions from Sections 1.1--1.3:
\begin{tabularx}{\linewidth}{@{}>{\bfseries}l X@{}}
$\Delta m =$ & \Blank[2.22in]\\[3pt]
$B =$ & \Blank[2.22in]\\[3pt]
$B/A =$ & \Blank[2.22in]
\end{tabularx}
Use ${}^{4}_{2}\mathrm{He}$ and the fixed mass data as a quick check.
\RuledLines{2}
\end{boardbox}

\switchcolumn

\HandoutSection{5}{Energy density of reactions}

\begin{fixedbox}{Plant-scale fixed inputs from the board notes}
\scriptsize
\begin{tabularx}{\linewidth}{@{}>{\bfseries}l X X@{}}
\toprule
 & nuclear plant & coal plant\\
\midrule
electric power & $1\unit{GW_e}$ & $1\unit{GW_e}$\\
efficiency & $0.30$ & $0.40$\\
material basis & $100\unit{t}$ $\mathrm{UO_2}$ over $5\unit{y}$ & $10{,}000\unit{t}$ per day\\
\bottomrule
\end{tabularx}
\end{fixedbox}

\begin{boardbox}{Energy per unit material mass}
Use $E=P\,\Delta t$ and account for thermal efficiency:
\[
\left.\frac{E}{m}\right|_{\mathrm{nuc}}
\approx
\frac{1\times10^9\unit{W_e}}{0.30}
\frac{5(365)(24)\unit{h}}{100\times10^3\unit{kg}}
=\RevealBlank[0.78in]{1.5 GWh/kg}
\]
\[
\left.\frac{E}{m}\right|_{\mathrm{coal}}
\approx
\frac{1\times10^9\unit{W_e}}{0.40}
\frac{24\unit{h}}{10{,}000\times10^3\unit{kg}}
=\RevealBlank[0.78in]{6 kWh/kg}
\]
\textbf{Ratio and interpretation:} \Blank[1.90in]
\end{boardbox}

\begin{boardbox}{Reaction-scale mass conversion}
\[
\left.\frac{\Delta m}{m_0}\right|_{\mathrm{fission}}
\approx \frac{200}{236(931.5)}
=\Blank[0.78in]
\]
\[
\left.\frac{\Delta m}{m_0}\right|_{\mathrm{combustion}}
\approx \frac{4\times10^{-6}}{44(931.5)}
=\Blank[0.78in]
\]
How do these fractions explain the plant-scale comparison?
\RuledLines{1}
\end{boardbox}

\begin{takeawaybox}[Three related scales]
For pure matter,
\[
  \frac{E}{m}=c^2\approx\RevealBlank[1.45in]{25,000 GWh/kg}.
\]
From the board notes,
\[
  \left.\frac{P}{m}\right|_{\mathrm{nuc}}
  \approx\RevealBlank[0.92in]{33 kW/kg},
  \qquad
  \left.\frac{P}{m}\right|_{\mathrm{coal}}
  \approx\RevealBlank[0.92in]{250 W/kg}.
\]
Why is the power-density contrast smaller than the energy-density contrast?
\RuledLines{1}
\end{takeawaybox}

% Instructor cue: distinguish energy density from power density and close by
% asking where batteries lie on the same two axes.

\end{handoutcolumns}
\end{document}
